% TOP_PROBLEM: {"id":30007699,"problem_number":"LOCAL-30007699","title":"Open versus proprietary AI","category_id":21,"category":{"id":21,"name":"economics","display_name":"Economics","description":"Open economic theory statements, published research agendas, and proposed scoped specifications; question provenance and historical priority are qualified per record.","slug":"economics","order_index":21,"created_at":"2026-10-08T00:00:00Z"},"status":"provenance_pending","external_url":null,"legacy_ids":[],"published":false,"statement":"Estimate when open model access increases downstream innovation net of real costs, including startup spillovers and upstream investment responses.","economics":{"original_id":188,"field":"Innovation and AI economics","kind":"empirical","formalization_basis":"proposed_specification","source_status":"not_verified","priority_status":"not_established","priority_note":"No exact open-problem passage was verified. The supporting publication does not establish priority or unresolved status.","earliest_verified_reference":null,"current_reference":{"authors":["Joshua S. Gans"],"title":"Market Power in Artificial Intelligence","year":2026,"url":"https://doi.org/10.1146/annurev-economics-051624-061832","doi":"10.1146/annurev-economics-051624-061832","locator":"Primary survey abstract and indexed excerpt; exact open-model innovation paragraph not retrieved","evidence_summary":"Surveys AI market power and data-market mechanisms; the open-versus-proprietary comparison is a proposed specification."},"formalization_note":"The cited primary work supports the subject or supplies solutions in particular settings, but no exact open-question passage for this specification was verified. The mathematical target and restrictions are proposed here, with no canonical-conjecture or priority claim."},"statusQualification":"Provenance pending: an explicit published statement of this exact open question was not verified. This is a catalogue-authored candidate specification, not a certified open conjecture. The cited primary work supports the subject or supplies solutions in particular settings, but no exact open-question passage for this specification was verified. The mathematical target and restrictions are proposed here, with no canonical-conjecture or priority claim.","statusReviewedAt":"2026-10-08"}
% FIRST_POSTED: 2026-10-08
% ENTRY_KIND: statement_only
% !TeX program = lualatex
\documentclass[11pt]{article}
\usepackage{fontspec}
\usepackage[a4paper,margin=25mm]{geometry}
\usepackage{amsmath,amssymb,mathtools}
\usepackage{xurl}
\usepackage[hidelinks]{hyperref}
\newcommand{\sref}[2]{\hyperlink{source:#1:#2}{[S#2]}}
\setlength{\emergencystretch}{3em}
\begin{document}
% problemId:30007699
\section*{Open versus proprietary AI}\label{30007699}
\subsection{Definitions and mathematical statement}
Consider downstream AI startups in prespecified product markets over three years. At baseline, compare two access regimes for models of matched measured capability: \(O\), which permits downloading, modification and redistribution under a stated license, and \(P\), which provides a priced proprietary service. Compute access and initial support are recorded separately. For market \(m\), let \(N_m(r)\) be independently validated new-product value under regime \(r\in\{O,P\}\), and \(C_m(r)\) all real development, adaptation and deployment costs. Define the downstream innovation effect \(\tau_N=E[N_m(O)-N_m(P)]\) and net-value effect \(\tau_W=E[N_m(O)-C_m(O)-N_m(P)+C_m(P)]\). Expectations concern eligible markets, not all AI development. Clustered rollout, credible licensing changes or another specified design must address endogenous selection into regimes and model-quality differences. Measure knowledge spillovers, entrant survival and incumbent responses, and separately include the upstream developer's future training investment when evaluating a sustained regime. Determine which combinations of appropriability, modification needs and complementary resources make \(\tau_W>0\), and whether short-run downstream gains persist once upstream incentives adjust. Existing theories of open and proprietary models support the comparison; this particular empirical target and general-equilibrium extension are proposed specifications.

\par\textbf{Formulation and scope.} The cited primary work supports the subject or supplies solutions in particular settings, but no exact open-question passage for this specification was verified. The mathematical target and restrictions are proposed here, with no canonical-conjecture or priority claim.

\par\textbf{Open-question provenance.} An explicit published open-question statement was not verified. Sources below are supporting literature and must not be treated as a first listing.

\par\textbf{Historical priority.} No exact open-problem passage was verified. The supporting publication does not establish priority or unresolved status.

\subsection{Short English statement}
Estimate when open model access increases downstream innovation net of real costs, including startup spillovers and upstream investment responses.

\subsection{Sources}
\begin{itemize}\raggedright
\item[\textnormal{[S1]}]\hypertarget{source:30007699:1}{} Joshua S. Gans. 2026. Market Power in Artificial Intelligence. Primary survey abstract and indexed excerpt; exact open-model innovation paragraph not retrieved.
\url{https://doi.org/10.1146/annurev-economics-051624-061832}.
DOI: 10.1146/annurev-economics-051624-061832.
{\small\textit{Supporting or current literature; see evidence qualification. Surveys AI market power and data-market mechanisms; the open-versus-proprietary comparison is a proposed specification.}}
\end{itemize}
\end{document}
